% The abstract of the notes: Overview, Learning objectives, Prerequisites.
% The headings and the objectives are \ensv pairs; the prose is in onlyin
% blocks, English first, each with its Swedish partner.

\emph{\ensv{Overview}{Översikt}} % sync english=d4b1ea swedish=c21045
\begin{onlyin}{english} % sync english=f1746a swedish=5d9f62
We start from a word you have met before, \emph{interface}, and take it
apart.  Then we do one task, listing the files of a directory, in three
ways: in a file manager, at a command line and in a program.  What
changes between them is who is on the other side: a person who looks and
points, a person who types, or a program.  The command line turns out to
serve a person and a program at once, and that is why the lectures that
follow are about it.  We end with a task that the lecture
\emph{Scripts} solves, and with the places where you will meet a
command line.
\end{onlyin}
\begin{onlyin}{swedish}
Vi börjar med ett ord som du har mött förut, \emph{gränssnitt}, och tar
isär det.  Sedan gör vi en uppgift, att lista filerna i en katalog, på
tre sätt: i en filhanterare, på en kommandorad och i ett program.  Det
som ändras mellan dem är vem som finns på andra sidan: en människa som
tittar och pekar, en människa som skriver, eller ett program.
Kommandoraden visar sig tjäna en människa och ett program på samma gång,
och det är därför de föreläsningar som följer handlar om den.  Vi slutar
med en uppgift som föreläsningen \emph{Skript} löser, och med de platser
där du kommer att möta en kommandorad.
\end{onlyin}

\emph{\ensv{Learning objectives}{Lärandemål}} % sync english=8412b2 swedish=525913
\ensv{After this lecture you should be able to:}%
  {Efter den här föreläsningen ska du kunna:} % sync english=78f8fb swedish=7fba3d
% Canvas outcome 10450: control the computer environment via the command
% line (learning-outcomes.md).
\begin{restatable}{lo}{InterfacesLOKinds}\label{InterfacesLOKinds}%
  \ensv{explain what an interface is, and tell a graphical user
    interface, a command-line interface and an application programming
    interface apart by who is on each side of them.}%
    {förklara vad ett gränssnitt är, och skilja ett grafiskt
    användargränssnitt, ett kommandoradsgränssnitt och ett
    programmeringsgränssnitt åt efter vem som finns på var sida om dem.} % sync english=6f2639 swedish=efa85b
\end{restatable}
% 10450
\begin{restatable}{lo}{InterfacesLOBoth}\label{InterfacesLOBoth}%
  \ensv{explain why a command-line interface serves both a person and a
    program, and give tasks where that matters.}%
    {förklara varför ett kommandoradsgränssnitt tjänar både en människa
    och ett program, och ge exempel på uppgifter där det spelar roll.} % sync english=94753d swedish=1fdf4b
\end{restatable}
% 10450
\begin{restatable}{lo}{InterfacesLOWhere}\label{InterfacesLOWhere}%
  \ensv{name settings where you will meet a command line, such as remote
    servers, computing clusters and the processing of research data.}%
    {nämna sammanhang där du kommer att möta en kommandorad, som
    fjärrservrar, beräkningskluster och bearbetning av forskningsdata.} % sync english=fc7dfb swedish=914256
\end{restatable}

\emph{\ensv{Prerequisites}{Förkunskaper}} % sync english=865514 swedish=192934
\begin{onlyin}{english} % sync english=2811ae swedish=4d1d8c
None beyond having used a computer through its graphical interface.  To
try the commands at your own keyboard you need a terminal; the course's
guide to accessing the terminal covers Linux, macOS and Windows.
\end{onlyin}
\begin{onlyin}{swedish}
Inga utöver att ha använt en dator genom dess grafiska gränssnitt.  För
att pröva kommandona vid ditt eget tangentbord behöver du en terminal;
kursens guide till terminalen tar upp Linux, macOS och Windows.
\end{onlyin}
